The matter of discerning the correct audience for a work is certainly one of the more important obligations of an author.
Disappointing in this regard can cause obfuscation for the reader of what the work seeks to convey, and thus the work ultimately not conveying anything of importance at all.
Obviously, the author ought then to strive for clarity for his intended audience to the greatest extent possible.
Unclear for him, however, is who this intended audience is meant to be for a master's thesis in fluid mechanics.
The immediate, and frankly short-sighted answer is of course the examiner.
For they shall decide whether this lamentable author and candidate for master of science is also worthy of pursuing a doctoral degree of the same variety.
If the work is unclear, then worthy he certainly is not.
Upon further reflexion, the author may realize that writing to appease an examiner could not be the purpose of the thesis---he perhaps figures that the intended audience is a natural result of the purpose of the thesis work.
For whom, then, does the author seek to obtain a master's degree?
The answer is obvious: \emph{himself}.

His pursual of a master's degree is not for the intellectual satisfaction of himself, rather being an exercise in the acquisition of new, hitherto unknown knowledge at the brink of his field of study.
Presentation of this accrued knowledge is the purpose of the accompanying thesis.
